\documentclass[10pt,a4paper]{article}
\PassOptionsToPackage{hyphens}{url}
\usepackage{amsmath,amssymb,amsthm,graphicx,verbatim,fontspec,hyperref}
\setmainfont{lmroman10-regular.otf}
\emergencystretch=3em
\usepackage[a4paper,margin=20mm]{geometry}
% Retained preamble for the Stacks Project chapter snapshot in this directory.
% The chapter's own preamble.tex is a copy of the shared Stacks Project preamble; the block
% below is sliced MECHANICALLY, line for line, from it. Nothing was added, removed or
% redefined; the wrapper class, the xr-hyper/hyperref wiring and the externaldocument list
% are omitted because the assembler supplies the document class.
%
% lines taken from preamble.tex: 3, 5, 9, 11, 12, 156, 157, 158, ... (51 lines total)
%
% package coverage restored from the chapter's own preamble (2026-10-13): xy, xy-2cell

\usepackage[all]{xy}
\xyoption{2cell}
\UseAllTwocells

\newenvironment{reference}{\comment}{\endcomment}
\newenvironment{slogan}{\comment}{\endcomment}
\newenvironment{history}{\comment}{\endcomment}
\theoremstyle{plain}
\newtheorem{theorem}[subsection]{Theorem}
\newtheorem{proposition}[subsection]{Proposition}
\newtheorem{lemma}[subsection]{Lemma}
\theoremstyle{definition}
\newtheorem{definition}[subsection]{Definition}
\newtheorem{example}[subsection]{Example}
\newtheorem{exercise}[subsection]{Exercise}
\newtheorem{situation}[subsection]{Situation}
\theoremstyle{remark}
\newtheorem{remark}[subsection]{Remark}
\newtheorem{remarks}[subsection]{Remarks}
\numberwithin{equation}{subsection}
\def\lim{\mathop{\mathrm{lim}}\nolimits}
\def\colim{\mathop{\mathrm{colim}}\nolimits}
\def\Spec{\mathop{\mathrm{Spec}}}
\def\Hom{\mathop{\mathrm{Hom}}\nolimits}
\def\Ext{\mathop{\mathrm{Ext}}\nolimits}
\def\SheafHom{\mathop{\mathcal{H}\!\mathit{om}}\nolimits}
\def\SheafExt{\mathop{\mathcal{E}\!\mathit{xt}}\nolimits}
\def\Sch{\mathit{Sch}}
\def\Mor{\mathop{\mathrm{Mor}}\nolimits}
\def\Ob{\mathop{\mathrm{Ob}}\nolimits}
\def\Sh{\mathop{\mathit{Sh}}\nolimits}
\def\NL{\mathop{N\!L}\nolimits}
\def\CH{\mathop{\mathrm{CH}}\nolimits}
\def\proetale{{pro\text{-}\acute{e}tale}}
\def\etale{{\acute{e}tale}}
\def\QCoh{\mathit{QCoh}}
\def\Ker{\mathop{\mathrm{Ker}}}
\def\Im{\mathop{\mathrm{Im}}}
\def\Coker{\mathop{\mathrm{Coker}}}
\def\Coim{\mathop{\mathrm{Coim}}}
\DeclareMathSymbol{\boxtimes}{\mathbin}{AMSa}{"02}
\def\QCohstack{\mathcal{QC}\!\mathit{oh}}
\def\Cohstack{\mathcal{C}\!\mathit{oh}}
\def\Spacesstack{\mathcal{S}\!\mathit{paces}}
\def\Quotfunctor{\mathrm{Quot}}
\def\Hilbfunctor{\mathrm{Hilb}}
\def\Curvesstack{\mathcal{C}\!\mathit{urves}}
\def\Polarizedstack{\mathcal{P}\!\mathit{olarized}}
\def\Complexesstack{\mathcal{C}\!\mathit{omplexes}}
\def\Pic{\mathop{\mathrm{Pic}}\nolimits}
\def\Picardstack{\mathcal{P}\!\mathit{ic}}
\def\Picardfunctor{\mathrm{Pic}}
\def\Deformationcategory{\mathcal{D}\!\mathit{ef}}
\title{The canonical triangulated structure on a compatible localization}
\author{The Stacks Project Authors}
\date{stacks-project:derived.tex@a04446e57ec1fbc252a871afcec7752fb2807b14; GNU FDL 1.2 or later}
\begin{document}
\maketitle
\noindent\textit{Editorial note.} Original theorem, complete author proof, surrounding source sections and shared bibliography retained. Cross-section/chapter references outside the collected sections display their exact original labels, without invented theorem numbers. No hypothesis or mathematical argument added or repaired.\par\medskip
\noindent\textit{Editorial note.} Standalone paragraph emergency stretch prevents an overfull inline-triangle paragraph; original mathematical tokens are unchanged.\par\medskip
\noindent\textit{Editorial note.} Original-label pointers may wrap at hyphens for standalone readability.\par\medskip
\noindent\textit{Allowed standard prerequisites (editorial).}\begin{itemize}
\item Additive categories, categories of fractions, triangulated categories and basic homological algebra. All selected theorem constructions and the surrounding local definitions and proofs are included. Standard external references retain their exact original labels and shared bibliography.. No primary localization or representing-object construction is assumed as background.
\end{itemize}\medskip
\section{Introduction}
\label{section-introduction}

\noindent
We first discuss triangulated categories and localization in triangulated
categories. Next, we prove that the homotopy category of complexes in an
additive category is a triangulated category. Once this is done we define
the derived category of an abelian category as the localization of the
homotopy category with respect to quasi-isomorphisms.
A good reference is Verdier's thesis \cite{Verdier}.



\section{Triangulated categories}
\label{section-triangulated-categories}

\noindent
Triangulated categories are a convenient tool to describe the type
of structure inherent in the derived category of an abelian category.
Some references are \cite{Verdier}, \cite{KS}, and \cite{Neeman}.




\section{The definition of a triangulated category}
\label{section-triangulated-definitions}

\noindent
In this section we collect most of the definitions concerning triangulated
and pre-triangulated categories.

\begin{definition}
\label{definition-triangle}
Let $\mathcal{D}$ be an additive category. Let
$[1] : \mathcal{D} \to \mathcal{D}$, $E \mapsto E[1]$
be an additive functor which is an auto-equivalence of $\mathcal{D}$.
\begin{enumerate}
\item A {\it triangle} is a sextuple
$(X, Y, Z, f, g, h)$ where $X, Y, Z \in \Ob(\mathcal{D})$ and
$f : X \to Y$, $g : Y \to Z$ and $h : Z \to X[1]$ are morphisms
of $\mathcal{D}$.
\item A {\it morphism of triangles}
$(X, Y, Z, f, g, h) \to (X', Y', Z', f', g', h')$
is given by morphisms $a : X \to X'$, $b : Y \to Y'$ and $c : Z \to Z'$
of $\mathcal{D}$ such that
$b \circ f = f' \circ a$, $c  \circ g = g' \circ b$ and
$a[1] \circ h = h' \circ c$.
\end{enumerate}
\end{definition}

\noindent
A morphism of triangles is visualized by the following
commutative diagram
$$
\xymatrix{
X \ar[r] \ar[d]^a &
Y \ar[r] \ar[d]^b &
Z \ar[r] \ar[d]^c &
X[1] \ar[d]^{a[1]} \\
X' \ar[r] &
Y' \ar[r] &
Z' \ar[r] &
X'[1]
}
$$
In the setting of Definition \ref{definition-triangle}, we
write $[0] = \text{id}$, for $n > 0$ we denote $[n]$ the $n$-fold
composition of $[1]$, we choose a quasi-inverse $[-1]$ of $[1]$, and
we set $[-n]$ equal to the $n$-fold composition of $[-1]$. Then
$\{[n]\}_{n \in \mathbf{Z}}$ is a collection of additive auto-equivalences
of $\mathcal{D}$ indexed by $n \in \mathbf{Z}$ such that we are given
isomorphisms of functors $[n] \circ [m] \cong [n + m]$.

\medskip\noindent
Here is the definition of a triangulated category as given in
Verdier's thesis.

\begin{definition}
\label{definition-triangulated-category}
A {\it triangulated category} consists of a triple
$(\mathcal{D}, \{[n]\}_{n\in \mathbf{Z}}, \mathcal{T})$
where
\begin{enumerate}
\item $\mathcal{D}$ is an additive category,
\item $[1] : \mathcal{D} \to \mathcal{D}$, $E \mapsto E[1]$
is an additive auto-equivalence and $[n]$ for $n \in \mathbf{Z}$ is
as discussed above, and
\item $\mathcal{T}$ is a set of triangles (Definition \ref{definition-triangle})
called the {\it distinguished triangles}
\end{enumerate}
subject to the following conditions
\begin{enumerate}
\item[TR1] Any triangle isomorphic to a distinguished triangle is
a distinguished triangle. Any triangle of the form
$(X, X, 0, \text{id}, 0, 0)$ is distinguished.
For any morphism $f : X \to Y$ of $\mathcal{D}$ there exists a
distinguished triangle of the form $(X, Y, Z, f, g, h)$.
\item[TR2] The triangle $(X, Y, Z, f, g, h)$ is distinguished
if and only if the triangle $(Y, Z, X[1], g, h, -f[1])$ is.
\item[TR3] Given a solid diagram
$$
\xymatrix{
X \ar[r]^f \ar[d]^a &
Y \ar[r]^g \ar[d]^b &
Z \ar[r]^h \ar@{-->}[d] &
X[1] \ar[d]^{a[1]} \\
X' \ar[r]^{f'} &
Y' \ar[r]^{g'} &
Z' \ar[r]^{h'} &
X'[1]
}
$$
whose rows are distinguished triangles and which satisfies
$b \circ f = f' \circ a$, there exists a morphism
$c : Z \to Z'$ such that $(a, b, c)$ is a morphism of triangles.
\item[TR4] Given objects $X$, $Y$, $Z$ of $\mathcal{D}$, and morphisms
$f : X \to Y$, $g : Y \to Z$, and distinguished triangles
$(X, Y, Q_1, f, p_1, d_1)$,
$(X, Z, Q_2, g \circ f, p_2, d_2)$,
and
$(Y, Z, Q_3, g, p_3, d_3)$,
there exist
morphisms $a : Q_1 \to Q_2$ and $b : Q_2 \to Q_3$ such
that
\begin{enumerate}
\item $(Q_1, Q_2, Q_3, a, b, p_1[1] \circ d_3)$ is a
distinguished triangle,
\item the triple $(\text{id}_X, g, a)$ is
a morphism of triangles
$(X, Y, Q_1, f, p_1, d_1) \to (X, Z, Q_2, g \circ f, p_2, d_2)$, and
\item the triple $(f, \text{id}_Z, b)$ is a morphism of triangles
$(X, Z, Q_2, g \circ f, p_2, d_2) \to (Y, Z, Q_3, g, p_3, d_3)$.
\end{enumerate}
\end{enumerate}
We will call $(\mathcal{D}, [\ ], \mathcal{T})$ a
{\it pre-triangulated category} if TR1, TR2 and TR3
hold.\footnote{We use $[\ ]$ as an abbreviation for the
family $\{[n]\}_{n\in \mathbf{Z}}$.}
\end{definition}

\noindent
The explanation of TR4 is that if you think of $Q_1$ as
$Y/X$, $Q_2$ as $Z/X$ and $Q_3$ as $Z/Y$, then TR4(a) expresses
the isomorphism $(Z/X)/(Y/X) \cong Z/Y$ and TR4(b) and TR4(c)
express that we can compare the triangles $X \to Y \to Q_1 \to X[1]$
etc with morphisms of triangles. For a more precise reformulation
of this idea see the proof of Lemma \textnormal{[original label: \nolinkurl{lemma-two-split-injections}]}.

\medskip\noindent
The sign in TR2 means that if $(X, Y, Z, f, g, h)$ is a distinguished triangle
then in the long sequence
\begin{equation}
\label{equation-rotate}
\ldots \to
Z[-1] \xrightarrow{-h[-1]}
X \xrightarrow{f}
Y \xrightarrow{g}
Z \xrightarrow{h}
X[1] \xrightarrow{-f[1]}
Y[1] \xrightarrow{-g[1]}
Z[1] \to \ldots
\end{equation}
each four term sequence gives a distinguished triangle.

\medskip\noindent
As usual we abuse notation and we simply speak of a (pre-)triangulated
category $\mathcal{D}$ without explicitly introducing notation for the
additional data. The notion of a pre-triangulated category is
useful in finding statements equivalent to TR4.

\medskip\noindent
We have the following definition of a triangulated functor.

\begin{definition}
\label{definition-exact-functor-triangulated-categories}
Let $\mathcal{D}$, $\mathcal{D}'$ be pre-triangulated
categories. An {\it exact functor}, or a {\it triangulated functor}
from $\mathcal{D}$ to $\mathcal{D}'$ is a functor
$F : \mathcal{D} \to \mathcal{D}'$ together
with given functorial isomorphisms $\xi_X : F(X[1]) \to F(X)[1]$
such that for every distinguished triangle
$(X, Y, Z, f, g, h)$ of $\mathcal{D}$ the triangle
$(F(X), F(Y), F(Z), F(f), F(g), \xi_X \circ F(h))$
is a distinguished triangle of $\mathcal{D}'$.
\end{definition}

\noindent
An exact functor is additive, see
Lemma \ref{lemma-exact-functor-additive}.
When we say two triangulated categories are equivalent we mean that
they are equivalent in the $2$-category of triangulated categories.
A $2$-morphism $a : (F, \xi) \to (F', \xi')$ in this $2$-category is
simply a transformation of functors $a : F \to F'$ which is compatible
with $\xi$ and $\xi'$, i.e.,
$$
\xymatrix{
F \circ [1] \ar[r]_\xi \ar[d]_{a \star 1} & [1] \circ F \ar[d]^{1 \star a} \\
F' \circ [1] \ar[r]^{\xi'} & [1] \circ F'
}
$$
commutes.

\begin{definition}
\label{definition-triangulated-subcategory}
Let $(\mathcal{D}, [\ ], \mathcal{T})$ be a pre-triangulated category.
A {\it pre-triangulated subcategory}\footnote{This definition may be
nonstandard. If $\mathcal{D}'$ is a full subcategory then $\mathcal{T}'$
is the intersection of the set of triangles in $\mathcal{D}'$ with
$\mathcal{T}$, see
Lemma \ref{lemma-triangulated-subcategory}.
In this case we drop $\mathcal{T}'$ from the notation.}
is a pair $(\mathcal{D}', \mathcal{T}')$ such that
\begin{enumerate}
\item $\mathcal{D}'$ is an additive subcategory of $\mathcal{D}$
which is preserved under $[1]$ and such that
$[1] : \mathcal{D}' \to \mathcal{D}'$ is an auto-equivalence,
\item $\mathcal{T}' \subset \mathcal{T}$ is a subset such that for every
$(X, Y, Z, f, g, h) \in \mathcal{T}'$ we have
$X, Y, Z \in \Ob(\mathcal{D}')$ and
$f, g, h \in \text{Arrows}(\mathcal{D}')$, and
\item $(\mathcal{D}', [\ ], \mathcal{T}')$ is a pre-triangulated
category.
\end{enumerate}
If $\mathcal{D}$ is a triangulated category, then we say
$(\mathcal{D}', \mathcal{T}')$ is a {\it triangulated subcategory} if
it is a pre-triangulated subcategory and
$(\mathcal{D}', [\ ], \mathcal{T}')$ is a triangulated category.
\end{definition}

\noindent
In this situation the inclusion functor
$\mathcal{D}' \to \mathcal{D}$ is an exact functor
with $\xi_X : X[1] \to X[1]$ given by the identity on $X[1]$.

\medskip\noindent
We will see in
Lemma \ref{lemma-composition-zero}
that for a distinguished triangle $(X, Y, Z, f, g, h)$
in a pre-triangulated category the composition $g \circ f : X \to Z$ is zero.
Thus the sequence (\ref{equation-rotate}) is a complex.
A homological functor is one that turns this complex into a long
exact sequence.

\begin{definition}
\label{definition-homological}
Let $\mathcal{D}$ be a pre-triangulated category.
Let $\mathcal{A}$ be an abelian category.
An additive functor $H : \mathcal{D} \to \mathcal{A}$ is called
{\it homological} if for every distinguished triangle
$(X, Y, Z, f, g, h)$ the sequence
$$
H(X) \to H(Y) \to H(Z)
$$
is exact in the abelian category $\mathcal{A}$. An additive functor
$H : \mathcal{D}^{opp} \to \mathcal{A}$ is called {\it cohomological}
if the corresponding functor $\mathcal{D} \to \mathcal{A}^{opp}$ is
homological.
\end{definition}

\noindent
If $H : \mathcal{D} \to \mathcal{A}$ is a homological functor
we often write $H^n(X) = H(X[n])$ so that $H(X) = H^0(X)$.
Our discussion of TR2 above implies that a distinguished triangle
$(X, Y, Z, f, g, h)$ determines a long exact sequence
\begin{equation}
\label{equation-long-exact-cohomology-sequence}
\xymatrix@C=3pc{
H^{-1}(Z) \ar[r]^{H(h[-1])} &
H^0(X) \ar[r]^{H(f)} &
H^0(Y) \ar[r]^{H(g)} &
H^0(Z) \ar[r]^{H(h)} &
H^1(X)
}
\end{equation}
This will be called the {\it long exact sequence} associated to the
distinguished triangle and the homological functor. As indicated
we will not use any signs for the morphisms in the long exact
sequence. This has the side effect that maps in the long exact sequence
associated to the rotation (TR2) of a distinguished triangle differ
from the maps in the sequence above by some signs.

\begin{definition}
\label{definition-delta-functor}
Let $\mathcal{A}$ be an abelian category.
Let $\mathcal{D}$ be a triangulated category.
A {\it $\delta$-functor from $\mathcal{A}$ to $\mathcal{D}$} is
given by a functor $G : \mathcal{A} \to \mathcal{D}$ and
a rule which assigns to every short exact sequence
$$
0 \to A \xrightarrow{a} B \xrightarrow{b} C \to 0
$$
a morphism $\delta = \delta_{A \to B \to C} : G(C) \to G(A)[1]$
such that
\begin{enumerate}
\item the triangle
$(G(A), G(B), G(C), G(a), G(b), \delta_{A \to B \to C})$
is a distinguished triangle of $\mathcal{D}$
for any short exact sequence as above, and
\item for every morphism $(A \to B \to C) \to (A' \to B' \to C')$
of short exact sequences the diagram
$$
\xymatrix{
G(C) \ar[d] \ar[rr]_{\delta_{A \to B \to C}} & &
G(A)[1] \ar[d] \\
G(C') \ar[rr]^{\delta_{A' \to B' \to C'}} & &
G(A')[1]
}
$$
is commutative.
\end{enumerate}
In this situation we call
$(G(A), G(B), G(C), G(a), G(b), \delta_{A \to B \to C})$
the {\it image of the short exact sequence under the
given $\delta$-functor}.
\end{definition}

\noindent
Note how a $\delta$-functor comes equipped with additional structure.
Strictly speaking it does not make sense to say that a given
functor $\mathcal{A} \to \mathcal{D}$ is a $\delta$-functor, but we
will often do so anyway.












\section{Elementary results on triangulated categories}
\label{section-elementary-results}

\noindent
Most of the results in this section are proved for pre-triangulated categories
and a fortiori hold in any triangulated category.

\begin{lemma}
\label{lemma-composition-zero}
Let $\mathcal{D}$ be a pre-triangulated category.
Let $(X, Y, Z, f, g, h)$ be a distinguished triangle.
Then $g \circ f = 0$,
$h \circ g = 0$ and $f[1] \circ h = 0$.
\end{lemma}

\begin{proof}
By TR1 we know $(X, X, 0, 1, 0, 0)$ is a distinguished triangle.
Apply TR3 to
$$
\xymatrix{
X \ar[r] \ar[d]^1 &
X \ar[r] \ar[d]^f &
0 \ar[r] \ar@{-->}[d] &
X[1] \ar[d]^{1[1]} \\
X \ar[r]^f &
Y \ar[r]^g &
Z \ar[r]^h &
X[1]
}
$$
Of course the dotted arrow is the zero map. Hence the commutativity of
the diagram implies that $g \circ f = 0$. For the other cases
rotate the triangle, i.e., apply TR2.
\end{proof}

\begin{lemma}
\label{lemma-representable-homological}
Let $\mathcal{D}$ be a pre-triangulated category.
For any object $W$ of $\mathcal{D}$ the functor
$\Hom_\mathcal{D}(W, -)$ is homological, and the functor
$\Hom_\mathcal{D}(-, W)$ is cohomological.
\end{lemma}

\begin{proof}
Consider a distinguished triangle $(X, Y, Z, f, g, h)$.
We have already seen that $g \circ f = 0$, see
Lemma \ref{lemma-composition-zero}.
Suppose $a : W \to Y$ is a morphism such that $g \circ a = 0$.
Then we get a commutative diagram
$$
\xymatrix{
W \ar[r]_1 \ar@{..>}[d]^b &
W \ar[r] \ar[d]^a &
0 \ar[r] \ar[d]^0 &
W[1] \ar@{..>}[d]^{b[1]} \\
X \ar[r] & Y \ar[r] & Z \ar[r] & X[1]
}
$$
Both rows are distinguished triangles (use TR1 for the top row).
Hence we can fill the dotted arrow $b$ (first rotate using TR2,
then apply TR3, and then rotate back). This proves the lemma.
\end{proof}

\begin{lemma}
\label{lemma-third-isomorphism-triangle}
Let $\mathcal{D}$ be a pre-triangulated category.
Let
$$
(a, b, c) : (X, Y, Z, f, g, h) \to (X', Y', Z', f', g', h')
$$
be a morphism of distinguished triangles. If two among $a, b, c$
are isomorphisms so is the third.
\end{lemma}

\begin{proof}
Assume that $a$ and $c$ are isomorphisms.
For any object $W$ of $\mathcal{D}$ write
$H_W( - ) = \Hom_\mathcal{D}(W, -)$.
Then we get a commutative diagram of abelian groups
$$
\xymatrix{
H_W(Z[-1]) \ar[r] \ar[d] &
H_W(X) \ar[r] \ar[d] &
H_W(Y) \ar[r] \ar[d] &
H_W(Z) \ar[r] \ar[d] &
H_W(X[1]) \ar[d] \\
H_W(Z'[-1]) \ar[r] &
H_W(X') \ar[r] &
H_W(Y') \ar[r] &
H_W(Z') \ar[r] &
H_W(X'[1])
}
$$
By assumption the right two and left two vertical arrows are bijective.
As $H_W$ is homological by
Lemma \ref{lemma-representable-homological}
and the five lemma
(Homology, Lemma \textnormal{[original label: \nolinkurl{homology-lemma-five-lemma}]})
it follows that the middle vertical arrow is an isomorphism.
Hence by Yoneda's lemma, see
Categories, Lemma \textnormal{[original label: \nolinkurl{categories-lemma-yoneda}]}
we see that $b$ is an isomorphism.
This implies the other cases by rotating (using TR2).
\end{proof}

\begin{remark}
\label{remark-special-triangles}
Let $\mathcal{D}$ be an additive category with translation functors $[n]$
as in Definition \ref{definition-triangle}. Let us call a triangle
$(X, Y, Z, f, g, h)$ {\it special}\footnote{This is nonstandard notation.}
if for every object $W$ of $\mathcal{D}$
the long sequence of abelian groups
$$
\ldots \to
\Hom_\mathcal{D}(W, X) \to
\Hom_\mathcal{D}(W, Y) \to
\Hom_\mathcal{D}(W, Z) \to
\Hom_\mathcal{D}(W, X[1]) \to \ldots
$$
is exact.  The proof of Lemma \ref{lemma-third-isomorphism-triangle}
shows that if
$$
(a, b, c) : (X, Y, Z, f, g, h) \to (X', Y', Z', f', g', h')
$$
is a morphism of special triangles and if two among $a, b, c$
are isomorphisms so is the third. There is a dual statement for
{\it co-special} triangles, i.e., triangles which turn into long
exact sequences on applying the functor $\Hom_\mathcal{D}(-, W)$.
Thus distinguished triangles are special and co-special, but in
general there are many more (co-)special triangles, than there are
distinguished triangles.
\end{remark}

\begin{lemma}
\label{lemma-third-map-square-zero}
Let $\mathcal{D}$ be a pre-triangulated category.
Let
$$
(0, b, 0), (0, b', 0) : (X, Y, Z, f, g, h) \to (X, Y, Z, f, g, h)
$$
be endomorphisms of a distinguished triangle. Then $bb' = 0$.
\end{lemma}

\begin{proof}
Picture
$$
\xymatrix{
X \ar[r] \ar[d]^0 &
Y \ar[r] \ar[d]^{b, b'} \ar@{..>}[ld]^\alpha &
Z \ar[r] \ar[d]^0 \ar@{..>}[ld]^\beta &
X[1] \ar[d]^0 \\
X \ar[r] & Y \ar[r] & Z \ar[r] & X[1]
}
$$
Applying
Lemma \ref{lemma-representable-homological}
we find dotted arrows $\alpha$ and $\beta$ such that
$b' = f \circ \alpha$ and $b = \beta \circ g$. Then
$bb' = \beta \circ g \circ f \circ \alpha = 0$
as $g \circ f = 0$ by
Lemma \ref{lemma-composition-zero}.
\end{proof}

\begin{lemma}
\label{lemma-third-map-idempotent}
Let $\mathcal{D}$ be a pre-triangulated category.
Let $(X, Y, Z, f, g, h)$ be a distinguished triangle.
If
$$
\xymatrix{
Z \ar[r]_h \ar[d]_c & X[1] \ar[d]^{a[1]} \\
Z \ar[r]^h & X[1]
}
$$
is commutative and $a^2 = a$, $c^2 = c$, then there exists a
morphism $b : Y \to Y$ with $b^2 = b$ such that
$(a, b, c)$ is an endomorphism of the triangle $(X, Y, Z, f, g, h)$.
\end{lemma}

\begin{proof}
By TR3 there exists a morphism $b'$ such that
$(a, b', c)$ is an endomorphism of $(X, Y, Z, f, g, h)$.
Then $(0, (b')^2 - b', 0)$ is also an endomorphism. By
Lemma \ref{lemma-third-map-square-zero}
we see that $(b')^2 - b'$ has square zero.
Set $b = b' - (2b' - 1)((b')^2 - b') = 3(b')^2 - 2(b')^3$.
A computation shows that $(a, b, c)$ is an endomorphism and
that $b^2 - b = (4(b')^2 - 4b' - 3)((b')^2 - b')^2 = 0$.
\end{proof}

\begin{lemma}
\label{lemma-cone-triangle-unique-isomorphism}
Let $\mathcal{D}$ be a pre-triangulated category.
Let $f : X \to Y$ be a morphism of $\mathcal{D}$.
There exists a distinguished triangle $(X, Y, Z, f, g, h)$ which
is unique up to (nonunique) isomorphism of triangles.
More precisely, given a second such distinguished triangle
$(X, Y, Z', f, g', h')$ there exists an isomorphism
$$
(1, 1, c) : (X, Y, Z, f, g, h) \longrightarrow (X, Y, Z', f, g', h')
$$
\end{lemma}

\begin{proof}
Existence by TR1. Uniqueness up to isomorphism by TR3 and
Lemma \ref{lemma-third-isomorphism-triangle}.
\end{proof}

\begin{lemma}
\label{lemma-uniqueness-third-arrow}
Let $\mathcal{D}$ be a pre-triangulated category. Let
$$
(a, b, c) : (X, Y, Z, f, g, h) \to (X', Y', Z', f', g', h')
$$
be a morphism of distinguished triangles. If one of the following
conditions holds
\begin{enumerate}
\item $\Hom(Y, X') = 0$,
\item $\Hom(Z, Y') = 0$,
\item $\Hom(X, X') = \Hom(Z, X') = 0$,
\item $\Hom(Z, X') = \Hom(Z, Z') = 0$, or
\item $\Hom(X[1], Z') = \Hom(Z, X') = 0$
\end{enumerate}
then $b$ is the unique morphism from $Y \to Y'$ such that
$(a, b, c)$ is a morphism of triangles.
\end{lemma}

\begin{proof}
If we have a second morphism of triangles $(a, b', c)$
then $(0, b - b', 0)$ is a morphism of triangles. Hence we
have to show: the only morphism $b : Y \to Y'$ such that
$X \to Y \to Y'$ and $Y \to Y' \to Z'$ are zero is $0$.
We will use Lemma \ref{lemma-representable-homological}
without further mention. In particular, condition (3) implies (1).
Given condition (1) if the composition $g' \circ b : Y \to Y' \to Z'$
is zero, then $b$ lifts to a morphism $Y \to X'$ which has to be zero.
This proves (1).

\medskip\noindent
The proof of (2) and (4) are dual to this argument.

\medskip\noindent
Assume (5). Consider the diagram
$$
\xymatrix{
X \ar[r]_f \ar[d]^0 &
Y \ar[r]_g \ar[d]^b &
Z \ar[r]_h \ar[d]^0 \ar@{..>}[ld]^\epsilon &
X[1] \ar[d]^0 \\
X' \ar[r]^{f'} &
Y' \ar[r]^{g'} &
Z' \ar[r]^{h'} &
X'[1]
}
$$
We may choose $\epsilon$ such that $b = \epsilon \circ g$.
Then $g' \circ \epsilon \circ g = 0$ which implies that
$g' \circ \epsilon = \delta \circ h$ for some
$\delta \in \Hom(X[1], Z')$. Since $\Hom(X[1], Z') = 0$
we conclude that $g' \circ \epsilon = 0$. Hence
$\epsilon = f' \circ \gamma$ for some $\gamma \in \Hom(Z, X')$.
Since $\Hom(Z, X') = 0$ we conclude that $\epsilon = 0$
and hence $b = 0$ as desired.
\end{proof}

\begin{lemma}
\label{lemma-third-object-zero}
Let $\mathcal{D}$ be a pre-triangulated category.
Let $f : X \to Y$ be a morphism of $\mathcal{D}$.
The following are equivalent
\begin{enumerate}
\item $f$ is an isomorphism,
\item $(X, Y, 0, f, 0, 0)$ is a distinguished triangle, and
\item for any distinguished triangle $(X, Y, Z, f, g, h)$ we have $Z = 0$.
\end{enumerate}
\end{lemma}

\begin{proof}
By TR1 the triangle $(X, X, 0, 1, 0, 0)$ is distinguished.
Let $(X, Y, Z, f, g, h)$ be a distinguished triangle.
By TR3 there is a map of distinguished triangles
$(1, f, 0) : (X, X, 0) \to (X, Y, Z)$.
If $f$ is an isomorphism, then $(1, f, 0)$ is an isomorphism
of triangles by Lemma \ref{lemma-third-isomorphism-triangle}
and $Z = 0$. Conversely, if $Z = 0$, then $(1, f, 0)$ is an
isomorphism of triangles as well, hence $f$ is an isomorphism.
\end{proof}

\begin{lemma}
\label{lemma-direct-sum-triangles}
Let $\mathcal{D}$ be a pre-triangulated category.
Let $(X, Y, Z, f, g, h)$ and $(X', Y', Z', f', g', h')$ be triangles.
The following are equivalent
\begin{enumerate}
\item $(X \oplus X', Y \oplus Y', Z \oplus Z',
f \oplus f', g \oplus g', h \oplus h')$
is a distinguished triangle,
\item both $(X, Y, Z, f, g, h)$ and $(X', Y', Z', f', g', h')$ are
distinguished triangles.
\end{enumerate}
\end{lemma}

\begin{proof}
Assume (2). By TR1 we may choose a distinguished triangle
$(X \oplus X', Y \oplus Y', Q, f \oplus f', g'', h'')$.
By TR3 we can find morphisms of distinguished triangles
$(X, Y, Z, f, g, h) \to
(X \oplus X', Y \oplus Y', Q, f \oplus f', g'', h'')$
and
$(X', Y', Z', f', g', h') \to
(X \oplus X', Y \oplus Y', Q, f \oplus f', g'', h'')$.
Taking the direct sum of these morphisms
we obtain a morphism of triangles
$$
\xymatrix{
(X \oplus X', Y \oplus Y', Z \oplus Z',
f \oplus f', g \oplus g', h \oplus h')
\ar[d]^{(1, 1, c)} \\
(X \oplus X', Y \oplus Y', Q, f \oplus f', g'', h'').
}
$$
In the terminology of Remark \ref{remark-special-triangles}
this is a map of special triangles (because a direct sum of special
triangles is special) and we conclude
that $c$ is an isomorphism. Thus (1) holds.

\medskip\noindent
Assume (1). We will show that $(X, Y, Z, f, g, h)$ is a distinguished
triangle. First observe that $(X, Y, Z, f, g, h)$ is a special triangle
(terminology from Remark \ref{remark-special-triangles})
as a direct summand of the distinguished hence special
triangle $(X \oplus X', Y \oplus Y', Z \oplus Z',
f \oplus f', g \oplus g', h \oplus h')$. Using TR1 let
$(X, Y, Q, f, g'', h'')$ be a distinguished triangle. By TR3 there exists
a morphism of distinguished triangles
 $(X \oplus X', Y \oplus Y', Z \oplus Z',
f \oplus f', g \oplus g', h \oplus h') \to (X, Y, Q, f, g'', h'')$.
Composing this with the inclusion map we get a morphism of triangles
$$
(1, 1, c) :
(X, Y, Z, f, g, h)
\longrightarrow
(X, Y, Q, f, g'', h'')
$$
By Remark \ref{remark-special-triangles}
we find that $c$ is an isomorphism and we conclude
that (2) holds.
\end{proof}

\begin{lemma}
\label{lemma-split}
Let $\mathcal{D}$ be a pre-triangulated category.
Let $(X, Y, Z, f, g, h)$ be a distinguished triangle.
\begin{enumerate}
\item If $h = 0$, then there exists a right inverse $s : Z \to Y$ to $g$.
\item For any right inverse $s : Z \to Y$ of $g$ the map
$f \oplus s : X \oplus Z \to Y$ is an isomorphism.
\item For any objects $X', Z'$ of $\mathcal{D}$ the triangle
$(X', X' \oplus Z', Z', (1, 0), (0, 1), 0)$ is distinguished.
\end{enumerate}
\end{lemma}

\begin{proof}
To see (1) use that
$\Hom_\mathcal{D}(Z, Y) \to \Hom_\mathcal{D}(Z, Z) \to
\Hom_\mathcal{D}(Z, X[1])$
is exact by
Lemma \ref{lemma-representable-homological}.
By the same token, if $s$ is as in (2), then $h = 0$ and the sequence
$$
0 \to \Hom_\mathcal{D}(W, X) \to \Hom_\mathcal{D}(W, Y)
\to \Hom_\mathcal{D}(W, Z) \to 0
$$
is split exact (split by $s : Z \to Y$). Hence by Yoneda's lemma we
see that $X \oplus Z \to Y$ is an isomorphism. The last assertion follows
from TR1 and
Lemma \ref{lemma-direct-sum-triangles}.
\end{proof}

\begin{lemma}
\label{lemma-when-split}
Let $\mathcal{D}$ be a pre-triangulated category.
Let $f : X \to Y$ be a morphism of $\mathcal{D}$.
The following are equivalent
\begin{enumerate}
\item $f$ has a kernel,
\item $f$ has a cokernel,
\item $f$ is the isomorphic to a composition
$K \oplus Z \to Z \to Z \oplus Q$ of a projection and coprojection
for some objects $K, Z, Q$ of $\mathcal{D}$.
\end{enumerate}
\end{lemma}

\begin{proof}
Any morphism isomorphic to a map of the form
$X' \oplus Z \to Z \oplus Y'$ has both a kernel and a cokernel.
Hence (3) $\Rightarrow$ (1), (2).
Next we prove (1) $\Rightarrow$ (3).
Suppose first that $f : X \to Y$ is a monomorphism, i.e., its kernel is zero.
By TR1 there exists a distinguished triangle $(X, Y, Z, f, g, h)$.
By Lemma \ref{lemma-composition-zero} the composition
$f \circ h[-1] = 0$. As $f$ is a monomorphism we see that $h[-1] = 0$
and hence $h = 0$. Then
Lemma \ref{lemma-split}
implies that $Y = X \oplus Z$, i.e., we see that (3) holds.
Next, assume $f$ has a kernel $K$. As $K \to X$ is a monomorphism we
conclude $X = K \oplus X'$ and $f|_{X'} : X' \to Y$ is a monomorphism.
Hence $Y = X' \oplus Y'$ and we win.
The implication (2) $\Rightarrow$ (3) is dual to this.
\end{proof}

\begin{lemma}
\label{lemma-products-sums-shifts-triangles}
Let $\mathcal{D}$ be a pre-triangulated category.
Let $I$ be a set.
\begin{enumerate}
\item Let $X_i$, $i \in I$ be a family of objects of $\mathcal{D}$.
\begin{enumerate}
\item If $\prod X_i$ exists, then $(\prod X_i)[1] = \prod X_i[1]$.
\item If $\bigoplus X_i$ exists, then $(\bigoplus X_i)[1] = \bigoplus X_i[1]$.
\end{enumerate}
\item Let $X_i \to Y_i \to Z_i \to X_i[1]$ be a family of distinguished
triangles of $\mathcal{D}$.
\begin{enumerate}
\item If $\prod X_i$, $\prod Y_i$, $\prod Z_i$ exist, then
$\prod X_i \to \prod Y_i \to \prod Z_i \to \prod X_i[1]$
is a distinguished triangle.
\item If $\bigoplus X_i$, $\bigoplus Y_i$,
$\bigoplus Z_i$ exist, then
$\bigoplus X_i \to \bigoplus Y_i \to \bigoplus Z_i \to \bigoplus X_i[1]$
is a distinguished triangle.
\end{enumerate}
\end{enumerate}
\end{lemma}

\begin{proof}
Part (1) is true because $[1]$ is an autoequivalence of $\mathcal{D}$
and because direct sums and products are defined in terms of the
category structure. Let us prove (2)(a). Choose a distinguished triangle
$\prod X_i \to \prod Y_i \to Z \to \prod X_i[1]$. For each $j$ we can
use TR3 to choose a morphism $p_j : Z \to Z_j$
fitting into a morphism of distinguished
triangles with the projection maps $\prod X_i \to X_j$ and $\prod Y_i \to Y_j$.
Using the definition of products we obtain a map
$\prod p_i : Z \to \prod Z_i$ fitting into a morphism
of triangles from the distinguished triangle to the triangle
made out of the products. Observe that the ``product'' triangle
$\prod X_i \to \prod Y_i \to \prod Z_i \to \prod X_i[1]$
is special in the terminology of Remark \ref{remark-special-triangles}
because products of exact sequences of abelian groups are exact.
Hence Remark \ref{remark-special-triangles} shows that
the morphism of triangles is an isomorphism and we conclude by TR1.
The proof of (2)(b) is dual.
\end{proof}

\begin{lemma}
\label{lemma-projectors-have-images-triangulated}
Let $\mathcal{D}$ be a pre-triangulated category.
If $\mathcal{D}$ has countable products, then $\mathcal{D}$
is Karoubian.
If $\mathcal{D}$ has countable coproducts, then $\mathcal{D}$
is Karoubian.
\end{lemma}

\begin{proof}
Assume $\mathcal{D}$ has countable products. By
Homology, Lemma \textnormal{[original label: \nolinkurl{homology-lemma-projectors-have-images}]}
it suffices to check that morphisms which have a right inverse have kernels.
Any morphism which has a right inverse is an epimorphism, hence
has a kernel by
Lemma \ref{lemma-when-split}.
The second statement is dual to the first.
\end{proof}

\noindent
The following lemma makes it slightly easier to prove that a
pre-triangulated category is triangulated.

\begin{lemma}
\label{lemma-easier-axiom-four}
Let $\mathcal{D}$ be a pre-triangulated category.
In order to prove TR4 it suffices to show that given
any pair of composable morphisms
$f : X \to Y$ and $g : Y \to Z$ there exist
\begin{enumerate}
\item isomorphisms $i : X' \to X$, $j : Y' \to Y$ and
$k : Z' \to Z$, and then setting $f' = j^{-1}fi : X' \to Y'$ and
$g' = k^{-1}gj : Y' \to Z'$ there exist
\item distinguished triangles
$(X', Y', Q_1, f', p_1, d_1)$,
$(X', Z', Q_2, g' \circ f', p_2, d_2)$
and
$(Y', Z', Q_3, g', p_3, d_3)$,
such that the assertion of TR4 holds.
\end{enumerate}
\end{lemma}

\begin{proof}
The replacement of $X, Y, Z$ by $X', Y', Z'$ is harmless by our
definition of distinguished triangles and their isomorphisms.
The lemma follows from the fact that the distinguished triangles
$(X', Y', Q_1, f', p_1, d_1)$,
$(X', Z', Q_2, g' \circ f', p_2, d_2)$
and
$(Y', Z', Q_3, g', p_3, d_3)$
are unique up to isomorphism by
Lemma \ref{lemma-cone-triangle-unique-isomorphism}.
\end{proof}

\begin{lemma}
\label{lemma-triangulated-subcategory}
Let $\mathcal{D}$ be a pre-triangulated category.
Assume that $\mathcal{D}'$ is an additive full subcategory of $\mathcal{D}$.
The following are equivalent
\begin{enumerate}
\item there exists a set of triangles $\mathcal{T}'$ such that
$(\mathcal{D}', \mathcal{T}')$ is a pre-triangulated subcategory
of $\mathcal{D}$,
\item $\mathcal{D}'$ is preserved under $[1]$ and
$[1] : \mathcal{D}' \to \mathcal{D}'$ is an auto-equivalence and
given any morphism $f : X \to Y$ in $\mathcal{D}'$ there exists
a distinguished triangle $(X, Y, Z, f, g, h)$ in $\mathcal{D}$
such that $Z$ is isomorphic to an object of $\mathcal{D}'$.
\end{enumerate}
In this case $\mathcal{T}'$ as in (1) is the set of distinguished triangles
$(X, Y, Z, f, g, h)$ of $\mathcal{D}$ such that
$X, Y, Z \in \Ob(\mathcal{D}')$. Finally, if $\mathcal{D}$
is a triangulated category, then (1) and (2) are also equivalent to
\begin{enumerate}
\item[(3)] $\mathcal{D}'$ is a triangulated subcategory.
\end{enumerate}
\end{lemma}

\begin{proof}
Omitted.
\end{proof}

\begin{lemma}
\label{lemma-exact-functor-additive}
An exact functor of pre-triangulated categories is additive.
\end{lemma}

\begin{proof}
Let $F : \mathcal{D} \to \mathcal{D}'$ be an exact functor of
pre-triangulated categories. Since
$(0, 0, 0, 1_0, 1_0, 0)$ is a distinguished triangle of $\mathcal{D}$
the triangle
$$
(F(0), F(0), F(0), 1_{F(0)}, 1_{F(0)}, F(0))
$$
is distinguished in $\mathcal{D}'$.
This implies that $1_{F(0)} \circ 1_{F(0)}$ is zero, see
Lemma \ref{lemma-composition-zero}.
Hence $F(0)$ is the zero object of $\mathcal{D}'$. This also implies
that $F$ applied to any zero morphism is zero (since a morphism in
an additive category is zero if and only if it factors through the
zero object). Next, using that
$(X, X \oplus Y, Y, (1, 0), (0, 1), 0)$ is a distinguished triangle by
Lemma \ref{lemma-split} part (3), we see that
$(F(X), F(X \oplus Y), F(Y), F(1, 0), F(0, 1), 0)$ is
one too. This implies that the map $F(X) \oplus F(Y) \to F(X \oplus Y)$
is an isomorphism by Lemma \ref{lemma-split} part (2). To finish we apply
Homology, Lemma \textnormal{[original label: \nolinkurl{homology-lemma-additive-functor}]}.
\end{proof}

\begin{lemma}
\label{lemma-exact-equivalence}
Let $F : \mathcal{D} \to \mathcal{D}'$ be a fully faithful exact functor
of pre-triangulated categories. Then a triangle $(X, Y, Z, f, g, h)$
of $\mathcal{D}$ is distinguished if and only if
$(F(X), F(Y), F(Z), F(f), F(g), F(h))$ is distinguished in $\mathcal{D}'$.
\end{lemma}

\begin{proof}
The ``only if'' part is clear. Assume $(F(X), F(Y), F(Z))$ is
distinguished in $\mathcal{D}'$. Pick a distinguished triangle
$(X, Y, Z', f, g', h')$ in $\mathcal{D}$. By
Lemma \ref{lemma-cone-triangle-unique-isomorphism}
there exists an isomorphism of triangles
$$
(1, 1, c') : (F(X), F(Y), F(Z)) \longrightarrow (F(X), F(Y), F(Z')).
$$
Since $F$ is fully faithful, there exists a morphism $c : Z \to Z'$
such that $F(c) = c'$. Then $(1, 1, c)$ is an isomorphism between
$(X, Y, Z)$ and $(X, Y, Z')$. Hence $(X, Y, Z)$ is distinguished
by TR1.
\end{proof}

\begin{lemma}
\label{lemma-composition-exact}
Let $\mathcal{D}, \mathcal{D}', \mathcal{D}''$ be pre-triangulated categories.
Let $F : \mathcal{D} \to \mathcal{D}'$ and
$F' : \mathcal{D}' \to \mathcal{D}''$ be exact functors.
Then $F' \circ F$ is an exact functor.
\end{lemma}

\begin{proof}
Omitted.
\end{proof}

\begin{lemma}
\label{lemma-exact-compose-homological-functor}
Let $\mathcal{D}$ be a pre-triangulated category.
Let $\mathcal{A}$ be an abelian category.
Let $H : \mathcal{D} \to \mathcal{A}$ be a homological functor.
\begin{enumerate}
\item Let $\mathcal{D}'$ be a pre-triangulated category.
Let $F : \mathcal{D}' \to \mathcal{D}$ be an exact functor.
Then the composition $H \circ F$ is a homological functor as well.
\item Let $\mathcal{A}'$ be an abelian category. Let
$G : \mathcal{A} \to \mathcal{A}'$ be an exact functor.
Then $G \circ H$ is a homological functor as well.
\end{enumerate}
\end{lemma}

\begin{proof}
Omitted.
\end{proof}

\begin{lemma}
\label{lemma-exact-compose-delta-functor}
Let $\mathcal{D}$ be a triangulated category.
Let $\mathcal{A}$ be an abelian category.
Let $G : \mathcal{A} \to \mathcal{D}$ be a $\delta$-functor.
\begin{enumerate}
\item Let $\mathcal{D}'$ be a triangulated category.
Let $F : \mathcal{D} \to \mathcal{D}'$ be an exact functor.
Then the composition $F \circ G$ is a $\delta$-functor as well.
\item Let $\mathcal{A}'$ be an abelian category. Let
$H : \mathcal{A}' \to \mathcal{A}$ be an exact functor.
Then $G \circ H$ is a $\delta$-functor as well.
\end{enumerate}
\end{lemma}

\begin{proof}
Omitted.
\end{proof}

\begin{lemma}
\label{lemma-compose-delta-functor-homological}
Let $\mathcal{D}$ be a triangulated category.
Let $\mathcal{A}$ and $\mathcal{B}$ be abelian categories.
Let $G : \mathcal{A} \to \mathcal{D}$ be a $\delta$-functor.
Let $H : \mathcal{D} \to \mathcal{B}$ be a homological functor.
Assume that $H^{-1}(G(A)) = 0$ for all $A$ in $\mathcal{A}$.
Then the collection
$$
\{H^n \circ G, H^n(\delta_{A \to B \to C})\}_{n \geq 0}
$$
is a $\delta$-functor from $\mathcal{A} \to \mathcal{B}$, see
Homology, Definition \textnormal{[original label: \nolinkurl{homology-definition-cohomological-delta-functor}]}.
\end{lemma}

\begin{proof}
The notation signifies the following. If
$0 \to A \xrightarrow{a} B \xrightarrow{b} C \to 0$ is
a short exact sequence in $\mathcal{A}$, then
$$
\delta = \delta_{A \to B \to C} : G(C) \to G(A)[1]
$$
is a morphism in $\mathcal{D}$ such that
$(G(A), G(B), G(C), a, b, \delta)$ is
a distinguished triangle, see
Definition \ref{definition-delta-functor}.
Then $H^n(\delta) : H^n(G(C)) \to H^n(G(A)[1]) = H^{n + 1}(G(A))$
is clearly functorial in the short exact sequence.
Finally, the long exact cohomology sequence
(\ref{equation-long-exact-cohomology-sequence})
combined with the vanishing of $H^{-1}(G(C))$
gives a long exact sequence
$$
0 \to H^0(G(A)) \to H^0(G(B)) \to H^0(G(C))
\xrightarrow{H^0(\delta)} H^1(G(A)) \to \ldots
$$
in $\mathcal{B}$ as desired.
\end{proof}

\noindent
The proof of the following result uses TR4.

\begin{proposition}
\label{proposition-9}
Let $\mathcal{D}$ be a triangulated category. Any commutative diagram
$$
\xymatrix{
X \ar[r] \ar[d] & Y \ar[d] \\
X' \ar[r] & Y'
}
$$
can be extended to a diagram
$$
\xymatrix{
X \ar[r] \ar[d] & Y \ar[r] \ar[d] & Z \ar[r] \ar[d] & X[1] \ar[d] \\
X' \ar[r] \ar[d] & Y' \ar[r] \ar[d] & Z' \ar[r] \ar[d] & X'[1] \ar[d] \\
X'' \ar[r] \ar[d] & Y'' \ar[r] \ar[d] & Z'' \ar[r] \ar[d] & X''[1] \ar[d] \\
X[1] \ar[r] & Y[1] \ar[r] & Z[1] \ar[r] & X[2]
}
$$
where all the squares are commutative, except for the lower right square
which is anticommutative. Moreover, each of the rows and columns are
distinguished triangles. Finally, the morphisms on the bottom row
(resp.\ right column) are obtained from the morphisms of the top row
(resp.\ left column) by applying $[1]$.
\end{proposition}

\begin{proof}
During this proof we avoid writing the arrows in order to make the proof
legible. Choose distinguished triangles
$(X, Y, Z)$, $(X', Y', Z')$, $(X, X', X'')$, $(Y, Y', Y'')$, and
$(X, Y', A)$. Note that the morphism $X \to Y'$ is both equal
to the composition $X \to Y \to Y'$ and equal to the composition
$X \to X' \to Y'$. Hence, we can find morphisms
\begin{enumerate}
\item $a : Z \to A$ and $b : A \to Y''$, and
\item $a' : X'' \to A$ and $b' : A \to Z'$
\end{enumerate}
as in TR4. Denote $c : Y'' \to Z[1]$ the composition
$Y'' \to Y[1] \to Z[1]$ and denote $c' : Z' \to X''[1]$ the composition
$Z' \to X'[1] \to X''[1]$. The conclusion of our application TR4
are that
\begin{enumerate}
\item $(Z, A, Y'', a, b, c)$, $(X'', A, Z', a', b', c')$
are distinguished triangles,
\item $(X, Y, Z) \to (X, Y', A)$,
$(X, Y', A) \to (Y, Y', Y'')$,
$(X, X', X'') \to (X, Y', A)$,
$(X, Y', A) \to (X', Y', Z')$
are morphisms of triangles.
\end{enumerate}
First using that
$(X, X', X'') \to (X, Y', A)$ and $(X, Y', A) \to (Y, Y', Y'')$.
are morphisms of triangles we see the first of the diagrams
$$
\vcenter{
\xymatrix{
X' \ar[r] \ar[d] & Y' \ar[d] \\
X'' \ar[r]^{b \circ a'} \ar[d] & Y'' \ar[d] \\
X[1] \ar[r] & Y[1]
}
}
\quad\text{and}\quad
\vcenter{
\xymatrix{
Y \ar[r] \ar[d] & Z \ar[d]^{b' \circ a} \ar[r] & X[1] \ar[d] \\
Y' \ar[r] & Z' \ar[r] & X'[1]
}
}
$$
is commutative. The second is commutative too using that
$(X, Y, Z) \to (X, Y', A)$ and $(X, Y', A) \to (X', Y', Z')$ are morphisms
of triangles. At this point we choose a distinguished triangle
$(X'', Y'' , Z'')$ starting with the map $b \circ a' : X'' \to Y''$.

\medskip\noindent
Next we apply TR4 one more time to the morphisms
$X'' \to A \to Y''$ and the triangles
$(X'', A, Z', a', b', c')$,
$(X'', Y'', Z'')$, and
$(A, Y'', Z[1], b, c , -a[1])$ to get morphisms
$a'' : Z' \to Z''$ and $b'' : Z'' \to Z[1]$.
Then $(Z', Z'', Z[1], a'', b'', - b'[1] \circ a[1])$ is a distinguished
triangle, hence also $(Z, Z', Z'', -b' \circ a, a'', -b'')$
and hence also $(Z, Z', Z'', b' \circ a, a'', b'')$.
Moreover, $(X'', A, Z') \to (X'', Y'', Z'')$ and
$(X'', Y'', Z'') \to (A, Y'', Z[1], b, c , -a[1])$
are morphisms of triangles.
At this point we have defined all the distinguished triangles
and all the morphisms, and all that's left is to verify some
commutativity relations.

\medskip\noindent
To see that the middle square in the diagram commutes, note
that the arrow $Y' \to Z'$ factors as $Y' \to A \to Z'$
because $(X, Y', A) \to (X', Y', Z')$ is a morphism of triangles.
Similarly, the morphism $Y' \to Y''$ factors as
$Y' \to A \to Y''$ because $(X, Y', A) \to (Y, Y', Y'')$ is a
morphism of triangles. Hence the middle square commutes because
the square with sides $(A, Z', Z'', Y'')$ commutes as
$(X'', A, Z') \to (X'', Y'', Z'')$ is a morphism of triangles (by TR4).
The square with sides $(Y'', Z'', Y[1], Z[1])$ commutes
because $(X'', Y'', Z'') \to (A, Y'', Z[1], b, c , -a[1])$
is a morphism of triangles and $c : Y'' \to Z[1]$ is the composition
$Y'' \to Y[1] \to Z[1]$.
The square with sides $(Z', X'[1], X''[1], Z'')$ is commutative
because $(X'', A, Z') \to (X'', Y'', Z'')$ is a morphism of triangles
and $c' : Z' \to X''[1]$ is the composition $Z' \to X'[1] \to X''[1]$.
Finally, we have to show that the square with sides
$(Z'', X''[1], Z[1], X[2])$ anticommutes. This holds because
$(X'', Y'', Z'') \to (A, Y'', Z[1], b, c , -a[1])$
is a morphism of triangles and we're done.
\end{proof}








\section{Localization of triangulated categories}
\label{section-localization}

\noindent
In order to construct the derived category starting from the homotopy
category of complexes, we will use a localization process.

\begin{definition}
\label{definition-localization}
Let $\mathcal{D}$ be a pre-triangulated category. We say a multiplicative
system $S$ is {\it compatible with the triangulated structure} if
the following two conditions hold:
\begin{enumerate}
\item[MS5] For a morphism $f$ of $\mathcal{D}$ we have
$f \in S \Leftrightarrow f[1] \in S$\footnote{See Remark \ref{remark-MS5}.}.
\item[MS6] Given a solid commutative square
$$
\xymatrix{
X \ar[r] \ar[d]^s &
Y \ar[r] \ar[d]^{s'} &
Z \ar[r] \ar@{-->}[d] &
X[1] \ar[d]^{s[1]} \\
X' \ar[r] &
Y' \ar[r] &
Z' \ar[r] &
X'[1]
}
$$
whose rows are distinguished triangles with $s, s' \in S$
there exists a morphism $s'' : Z \to Z'$ in $S$ such that
$(s, s', s'')$ is a morphism of triangles.
\end{enumerate}
\end{definition}

\noindent
It turns out that these axioms are not independent of the
axioms defining multiplicative systems.

\begin{lemma}
\label{lemma-localization-conditions}
Let $\mathcal{D}$ be a pre-triangulated category.
Let $S \subset \text{Arrows}(\mathcal{D})$.
\begin{enumerate}
\item If $S$ contains all identities and MS6 holds
(Definition \ref{definition-localization}),
then every isomorphism of $\mathcal{D}$ is in $S$.
\item If MS1, MS5 (Categories, Definition
\textnormal{[original label: \nolinkurl{categories-definition-multiplicative-system}]}) and MS6 hold,
then MS2 holds.
\end{enumerate}
\end{lemma}

\begin{proof}
Assume $S$ contains all identities and MS6 holds.
Let $f : X \to Y$ be an isomorphism of $\mathcal{D}$.
Consider the diagram
$$
\xymatrix{
0 \ar[r] \ar[d]^1 &
X \ar[r]_1 \ar[d]^1 &
X \ar[r] \ar@{-->}[d] &
0[1] \ar[d]^{1[1]} \\
0 \ar[r] &
X \ar[r]^f &
Y \ar[r] &
0[1]
}
$$
The rows are distinguished triangles by Lemma \ref{lemma-third-object-zero}.
By MS6 we see that the dotted arrow exists and is in $S$, so $f$ is in $S$.

\medskip\noindent
Assume MS1, MS5, MS6.
Suppose that $f : X \to Y$ is a morphism of $\mathcal{D}$ and
$t : X \to X'$ an element of $S$. Choose a distinguished triangle
$(X, Y, Z, f, g, h)$. Next, choose a distinguished triangle
$(X', Y', Z, f', g', t[1] \circ h)$ (here we use TR1 and TR2).
By MS5, MS6 (and TR2 to rotate) we can find the dotted arrow
in the commutative diagram
$$
\xymatrix{
X \ar[r] \ar[d]^t &
Y \ar[r] \ar@{..>}[d]^{s'} &
Z \ar[r] \ar[d]^1 &
X[1] \ar[d]^{t[1]} \\
X' \ar[r] &
Y' \ar[r] &
Z \ar[r] &
X'[1]
}
$$
with moreover $s' \in S$. This proves LMS2. The proof of RMS2 is dual.
\end{proof}

\begin{remark}
\label{remark-MS5}
In the presence of MS1 and MS6, condition MS5 is equivalent to asking
$s[n] \in S$ for all $s \in S$ and $n \in \mathbf{Z}$. For example,
suppose MS5 holds, we have $s \in S$, and we want to show $s[-1] \in S$.
This isn't immediate because $s[-1][1]$ is not equal to $s$, only
isomorphic to $s$ as an arrow of $\mathcal{D}$. Still, this does imply that
$s[-1][1] = f \circ s \circ g$ for isomorphisms $f$, $g$. By
Lemma \ref{lemma-localization-conditions} (1) we find $f, g \in S$,
hence $s[-1][1] \in S$ by MS1, hence $s[-1] \in S$ by MS5.
We leave a complete proof to the reader as an exercise.
\end{remark}

\begin{lemma}
\label{lemma-triangle-functor-localize}
Let $F : \mathcal{D} \to \mathcal{D}'$ be an exact functor of
pre-triangulated categories.  Let
$$
S = \{f \in \text{Arrows}(\mathcal{D})
\mid F(f)\text{ is an isomorphism}\}
$$
Then $S$ is a saturated (see
Categories,
Definition \textnormal{[original label: \nolinkurl{categories-definition-saturated-multiplicative-system}]})
multiplicative system compatible with the
triangulated structure on $\mathcal{D}$.
\end{lemma}

\begin{proof}
We have to prove axioms MS1 -- MS6, see
Categories, Definitions \textnormal{[original label: \nolinkurl{categories-definition-multiplicative-system}]} and
\textnormal{[original label: \nolinkurl{categories-definition-saturated-multiplicative-system}]}
and
Definition \ref{definition-localization}.
MS1, MS4, and MS5 are direct from the definitions. MS6 follows from TR3 and
Lemma \ref{lemma-third-isomorphism-triangle}.
By
Lemma \ref{lemma-localization-conditions}
we conclude that MS2 holds. To finish the proof we have to show that
MS3 holds. To do this let $f, g : X \to Y$ be morphisms of $\mathcal{D}$,
and let $t : Z \to X$ be an element of $S$ such that $f \circ t = g \circ t$.
As $\mathcal{D}$ is additive this simply means that $a \circ t = 0$ with
$a = f - g$. Choose a distinguished triangle $(Z, X, Q, t, d, h)$ using TR1.
Since $a \circ t = 0$ we see by
Lemma \ref{lemma-representable-homological}
there exists a morphism $i : Q \to Y$ such that $i \circ d = a$.
Finally, using TR1 again we can choose a triangle
$(Q, Y, W, i, j, k)$. Here is a picture
$$
\xymatrix{
Z \ar[r]_t & X \ar[r]_d \ar[d]^1 & Q \ar[r] \ar[d]^i & Z[1] \\
& X \ar[r]_a & Y \ar[d]^j \\
& & W
}
$$
OK, and now we apply the functor $F$ to this diagram.
Since $t \in S$ we see that $F(Q) = 0$, see
Lemma \ref{lemma-third-object-zero}.
Hence $F(j)$ is an isomorphism by the same lemma, i.e., $j \in S$.
Finally, $j \circ a = j \circ i \circ d = 0$ as $j \circ i = 0$.
Thus $j \circ f = j \circ g$ and we see that LMS3 holds.
The proof of RMS3 is dual.
\end{proof}

\begin{lemma}
\label{lemma-homological-functor-localize}
Let $H : \mathcal{D} \to \mathcal{A}$ be a homological functor between a
pre-triangulated category and an abelian category. Let
$$
S = \{f \in \text{Arrows}(\mathcal{D})
\mid H^i(f)\text{ is an isomorphism for all }i \in \mathbf{Z}\}
$$
Then $S$ is a saturated (see
Categories,
Definition \textnormal{[original label: \nolinkurl{categories-definition-saturated-multiplicative-system}]})
multiplicative system compatible with the
triangulated structure on $\mathcal{D}$.
\end{lemma}

\begin{proof}
We have to prove axioms MS1 -- MS6, see
Categories, Definitions \textnormal{[original label: \nolinkurl{categories-definition-multiplicative-system}]} and
\textnormal{[original label: \nolinkurl{categories-definition-saturated-multiplicative-system}]}
and
Definition \ref{definition-localization}.
MS1, MS4, and MS5 are direct from the definitions.
MS6 follows from TR3 and the long exact cohomology sequence
(\ref{equation-long-exact-cohomology-sequence}).
By
Lemma \ref{lemma-localization-conditions}
we conclude that MS2 holds. To finish the proof we have to show that
MS3 holds. To do this let $f, g : X \to Y$ be morphisms of $\mathcal{D}$,
and let $t : Z \to X$ be an element of $S$ such that $f \circ t = g \circ t$.
As $\mathcal{D}$ is additive this simply means that $a \circ t = 0$ with
$a = f - g$. Choose a distinguished triangle $(Z, X, Q, t, u, v)$ using
TR1 and TR2. Since $a \circ t = 0$ we see by
Lemma \ref{lemma-representable-homological}
there exists a morphism $i : Q \to Y$ such that $i \circ g = a$.
Finally, using TR1 again we can choose a triangle
$(Q, Y, W, i, j, k)$. Here is a picture
$$
\xymatrix{
Z \ar[r]_t & X \ar[r]_u \ar[d]^1 & Q \ar[r]_v \ar[d]^i & Z[1] \\
& X \ar[r]_a & Y \ar[d]^j \\
& & W
}
$$
OK, and now we apply the functors $H^i$ to this diagram.
Since $t \in S$ we see that $H^i(Q) = 0$ by the long exact cohomology
sequence (\ref{equation-long-exact-cohomology-sequence}).
Hence $H^i(j)$ is an isomorphism for all $i$ by the same argument,
i.e., $j \in S$. Finally, $j \circ a = j \circ i \circ u = 0$ as
$j \circ i = 0$. Thus $j \circ f = j \circ g$ and we see that LMS3 holds.
The proof of RMS3 is dual.
\end{proof}


\begin{proposition}
\label{proposition-construct-localization}
Let $\mathcal{D}$ be a pre-triangulated category. Let $S$ be a multiplicative
system compatible with the triangulated structure.
Then there exists a unique structure of a pre-triangulated category on
$S^{-1}\mathcal{D}$ such that $[1] \circ Q = Q \circ [1]$ and
the localization functor $Q : \mathcal{D} \to S^{-1}\mathcal{D}$ is exact.
Moreover, if $\mathcal{D}$ is a triangulated category, so is
$S^{-1}\mathcal{D}$.
\end{proposition}



\begin{proof}
We have seen that $S^{-1}\mathcal{D}$ is an additive category
and that the localization functor $Q$ is additive in
Homology, Lemma \textnormal{[original label: \nolinkurl{homology-lemma-localization-additive}]}.
It follows from MS5 that there is a unique additive auto-equivalence
$[1] : S^{-1}\mathcal{D} \to S^{-1}\mathcal{D}$
such that $Q \circ [1] = [1] \circ Q$ (equality of functors);
we omit the details. We say a triangle of $S^{-1}\mathcal{D}$
is distinguished if it is isomorphic to the image of a
distinguished triangle under the localization functor $Q$.

\medskip\noindent
Proof of TR1. The only thing to prove here is that if
$a : Q(X) \to Q(Y)$ is a morphism of $S^{-1}\mathcal{D}$, then
$a$ fits into a distinguished triangle. Write $a = Q(s)^{-1} \circ Q(f)$ for
some $s : Y \to Y'$ in $S$ and $f : X \to Y'$. Choose a distinguished
triangle $(X, Y', Z, f, g, h)$ in $\mathcal{D}$. Then we see that
$(Q(X), Q(Y), Q(Z), a, Q(g) \circ Q(s), Q(h))$ is a distinguished triangle
of $S^{-1}\mathcal{D}$.

\medskip\noindent
Proof of TR2. This is immediate from the definitions.

\medskip\noindent
Proof of TR3. Note that the existence of the dotted arrow which is
required to exist may be proven after replacing the two triangles
by isomorphic triangles. Hence we may assume given distinguished
triangles $(X, Y, Z, f, g, h)$ and $(X', Y', Z', f', g', h')$ of
$\mathcal{D}$ and a commutative diagram
$$
\xymatrix{
Q(X) \ar[r]_{Q(f)} \ar[d]_a & Q(Y) \ar[d]^b \\
Q(X') \ar[r]^{Q(f')} & Q(Y')
}
$$
in $S^{-1}\mathcal{D}$. Now we apply
Categories, Lemma \textnormal{[original label: \nolinkurl{categories-lemma-left-localization-diagram}]}
to find a morphism $f'' : X'' \to Y''$ in $\mathcal{D}$ and a commutative
diagram
$$
\xymatrix{
X \ar[d]_f \ar[r]_k & X'' \ar[d]^{f''} & X' \ar[d]^{f'} \ar[l]^s \\
Y \ar[r]^l & Y'' & Y' \ar[l]_t
}
$$
in $\mathcal{D}$ with $s, t \in S$ and $a = s^{-1}k$, $b = t^{-1}l$.
At this point we can use TR3 for $\mathcal{D}$ and MS6 to find
a commutative diagram
$$
\xymatrix{
X \ar[r] \ar[d]^k &
Y \ar[r] \ar[d]^l &
Z \ar[r] \ar[d]^m &
X[1] \ar[d]^{g[1]} \\
X'' \ar[r] &
Y'' \ar[r] &
Z'' \ar[r] &
X''[1] \\
X' \ar[r] \ar[u]_s &
Y' \ar[r] \ar[u]_t &
Z' \ar[r] \ar[u]_r &
X'[1] \ar[u]_{s[1]}
}
$$
with $r \in S$. It follows that setting $c = Q(r)^{-1}Q(m)$ we obtain
the desired morphism of triangles
$$
\xymatrix{
(Q(X), Q(Y), Q(Z), Q(f), Q(g), Q(h))
\ar[d]^{(a, b, c)} \\
(Q(X'), Q(Y'), Q(Z'), Q(f'), Q(g'), Q(h'))
}
$$

\medskip\noindent
This proves the first statement of the lemma. If $\mathcal{D}$ is also
a triangulated category, then we still have to prove TR4 in order to show
that $S^{-1}\mathcal{D}$ is triangulated as well. To do this we reduce by
Lemma \ref{lemma-easier-axiom-four}
to the following statement: Given composable morphisms
$a : Q(X) \to Q(Y)$ and $b : Q(Y) \to Q(Z)$ we have to produce
an octahedron after possibly replacing $Q(X), Q(Y), Q(Z)$ by isomorphic
objects. To do this we may first replace $Y$ by an object such that
$a = Q(f)$ for some morphism $f : X \to Y$ in $\mathcal{D}$. (More precisely,
write $a = s^{-1}f$ with $s : Y \to Y'$ in $S$ and $f : X \to Y'$. Then
replace $Y$ by $Y'$.) After this we similarly replace $Z$ by an object such
that $b = Q(g)$ for some morphism $g : Y \to Z$. Now we can find
distinguished triangles $(X, Y, Q_1, f, p_1, d_1)$,
$(X, Z, Q_2, g \circ f, p_2, d_2)$, and
$(Y, Z, Q_3, g, p_3, d_3)$ in $\mathcal{D}$ (by TR1), and
morphisms $a : Q_1 \to Q_2$ and $b : Q_2 \to Q_3$ as in TR4.
Then it is immediately verified that applying the functor $Q$ to
all these data gives a corresponding structure in $S^{-1}\mathcal{D}$.
\end{proof}


\noindent
The universal property of the localization of a triangulated category
is as follows (we formulate this for pre-triangulated categories, hence
it holds a fortiori for triangulated categories).

\begin{lemma}
\label{lemma-universal-property-localization}
Let $\mathcal{D}$ be a pre-triangulated category. Let $S$ be a multiplicative
system compatible with the triangulated structure. Let
$Q : \mathcal{D} \to S^{-1}\mathcal{D}$ be the localization functor, see
Proposition \ref{proposition-construct-localization}.
\begin{enumerate}
\item If $H : \mathcal{D} \to \mathcal{A}$ is a homological functor into
an abelian category $\mathcal{A}$ such that $H(s)$ is an isomorphism for
all $s \in S$, then the unique factorization
$H' : S^{-1}\mathcal{D} \to \mathcal{A}$ such that $H = H' \circ Q$ (see
Categories, Lemma \textnormal{[original label: \nolinkurl{categories-lemma-properties-left-localization}]})
is a homological functor too.
\item If $F : \mathcal{D} \to \mathcal{D}'$ is an exact functor into
a pre-triangulated category $\mathcal{D}'$ such that $F(s)$ is an isomorphism
for all $s \in S$, then the unique factorization
$F' : S^{-1}\mathcal{D} \to \mathcal{D}'$ such that $F = F' \circ Q$ (see
Categories, Lemma \textnormal{[original label: \nolinkurl{categories-lemma-properties-left-localization}]})
is an exact functor too.
\end{enumerate}
\end{lemma}

\begin{proof}
This lemma proves itself. Details omitted.
\end{proof}

\begin{lemma}
\label{lemma-localization-subcategory}
Let $\mathcal{D}$ be a pre-triangulated category and let
$\mathcal{D}' \subset \mathcal{D}$ be a full, pre-triangulated subcategory.
Let $S$ be a saturated multiplicative system of $\mathcal{D}$
compatible with the triangulated structure. Assume that for each $X$
in $\mathcal{D}$ there exists an $s : X' \to X$ in $S$ such that $X'$ is
an object of $\mathcal{D}'$. Then
$S' = S \cap \text{Arrows}(\mathcal{D}')$ is a saturated
multiplicative system compatible with the triangulated structure and
the functor
$$
(S')^{-1}\mathcal{D}' \longrightarrow S^{-1}\mathcal{D}
$$
is an equivalence of pre-triangulated categories.
\end{lemma}

\begin{proof}
Consider the quotient functor $Q : \mathcal{D} \to S^{-1}\mathcal{D}$
of Proposition \ref{proposition-construct-localization}.
Since $S$ is saturated we have that a morphism $f : X \to Y$ is in
$S$ if and only if $Q(f)$ is invertible, see
Categories, Lemma \textnormal{[original label: \nolinkurl{categories-lemma-what-gets-inverted}]}.
Thus $S'$ is the collection of arrows which are turned into isomorphisms
by the composition $\mathcal{D}' \to \mathcal{D} \to S^{-1}\mathcal{D}$.
Hence $S'$ is is a saturated multiplicative system compatible with the
triangulated structure by Lemma \ref{lemma-triangle-functor-localize}.
By Lemma \ref{lemma-universal-property-localization} we obtain the
exact functor $(S')^{-1}\mathcal{D}' \to S^{-1}\mathcal{D}$ of pre-triangulated
categories. By assumption this functor is essentially surjective.
Let $X', Y'$ be objects of $\mathcal{D}'$. By Categories, Remark
\textnormal{[original label: \nolinkurl{categories-remark-right-localization-morphisms-colimit}]} we have
$$
\Mor_{S^{-1}\mathcal{D}}(X', Y') =
\colim_{s : X \to X'\text{ in }S} \Mor_\mathcal{D}(X, Y')
$$
Our assumption implies that for any $s : X \to X'$ in $S$ we can
find a morphism $s' : X'' \to X$ in $S$ with $X''$ in $\mathcal{D}'$.
Then $s \circ s' : X'' \to X'$ is in $S'$. Hence the colimit above is
equal to
$$
\colim_{s'' : X'' \to X'\text{ in }S'} \Mor_{\mathcal{D}'}(X'', Y') =
\Mor_{(S')^{-1}\mathcal{D}'}(X', Y')
$$
This proves our functor is also fully faithful and the proof is complete.
\end{proof}

\noindent
The following lemma describes the kernel (see
Definition \textnormal{[original label: \nolinkurl{definition-kernel-category}]})
of the localization functor.

\begin{lemma}
\label{lemma-kernel-localization}
Let $\mathcal{D}$ be a pre-triangulated category. Let $S$ be a multiplicative
system compatible with the triangulated structure. Let $Z$ be an object
of $\mathcal{D}$. The following are equivalent
\begin{enumerate}
\item $Q(Z) = 0$ in $S^{-1}\mathcal{D}$,
\item there exists $Z' \in \Ob(\mathcal{D})$ such that
$0 : Z \to Z'$ is an element of $S$,
\item there exists $Z' \in \Ob(\mathcal{D})$ such that
$0 : Z' \to Z$ is an element of $S$, and
\item there exists an object $Z'$ and a distinguished triangle
$(X, Y, Z \oplus Z', f, g, h)$ such that $f \in S$.
\end{enumerate}
If $S$ is saturated, then these are also equivalent to
\begin{enumerate}
\item[(5)] the morphism $0 \to Z$ is an element of $S$,
\item[(6)] the morphism $Z \to 0$ is an element of $S$,
\item[(7)] there exists a distinguished triangle $(X, Y, Z, f, g, h)$
such that $f \in S$.
\end{enumerate}
\end{lemma}

\begin{proof}
The equivalence of (1), (2), and (3) is
Homology, Lemma \textnormal{[original label: \nolinkurl{homology-lemma-kernel-localization}]}.
If (2) holds, then $(Z'[-1], Z'[-1] \oplus Z, Z, (1, 0), (0, 1), 0)$
is a distinguished triangle (see
Lemma \ref{lemma-split})
with ``$0 \in S$''. By rotating we conclude that (4) holds.
If $(X, Y, Z \oplus Z', f, g, h)$ is a distinguished triangle with $f \in S$
then $Q(f)$ is an isomorphism hence $Q(Z \oplus Z') = 0$ hence $Q(Z) = 0$.
Thus (1) -- (4) are all equivalent.

\medskip\noindent
Next, assume that $S$ is saturated. Note that each of (5), (6), (7)
implies one of the equivalent conditions (1) -- (4). Suppose that
$Q(Z) = 0$. Then $0 \to Z$ is a morphism of $\mathcal{D}$ which becomes
an isomorphism in $S^{-1}\mathcal{D}$. According to
Categories, Lemma \textnormal{[original label: \nolinkurl{categories-lemma-what-gets-inverted}]}
the fact that $S$ is saturated implies that $0 \to Z$ is in $S$.
Hence (1) $\Rightarrow$ (5). Dually (1) $\Rightarrow$ (6).
Finally, if $0 \to Z$ is in $S$, then the triangle
$(0, Z, Z, 0, \text{id}_Z, 0)$ is distinguished by TR1 and TR2 and
is a triangle as in (7).
\end{proof}

\begin{lemma}
\label{lemma-limit-triangles}
Let $\mathcal{D}$ be a pre-triangulated category.
Let $S$ be a saturated multiplicative system in $\mathcal{D}$
that is compatible with the triangulated structure.
Let $(X, Y, Z, f, g, h)$ be a distinguished triangle in $\mathcal{D}$.
Consider the category of morphisms of triangles
$$
\mathcal{I} =
\{(s, s', s'') : (X, Y, Z, f, g, h) \to (X', Y', Z', f', g', h')
\mid s, s', s'' \in S\}
$$
Then $\mathcal{I}$ is a filtered category and the functors
$\mathcal{I} \to X/S$, $\mathcal{I} \to Y/S$, and $\mathcal{I} \to Z/S$
are cofinal.
\end{lemma}

\begin{proof}
We strongly suggest the reader skip the proof of this lemma and instead
work it out on a napkin.

\medskip\noindent
The first remark is that using rotation of distinguished triangles (TR2)
gives an equivalence of categories between $\mathcal{I}$ and the
corresponding category for the distinguished triangle
$(Y, Z, X[1], g, h, -f[1])$. Using this we see for example that if
we prove the functor $\mathcal{I} \to X/S$ is cofinal, then
the same thing is true for the functors $\mathcal{I} \to Y/S$ and
$\mathcal{I} \to Z/S$.

\medskip\noindent
Note that if $s : X \to X'$ is a morphism of $S$, then using
MS2 we can find $s' : Y \to Y'$ and $f' : X' \to Y'$ such that
$f' \circ s = s' \circ f$, whereupon we can use MS6 to complete
this into an object of $\mathcal{I}$. Hence the functor
$\mathcal{I} \to X/S$ is surjective on objects. Using rotation as above
this implies the same thing is true for the functors
$\mathcal{I} \to Y/S$ and $\mathcal{I} \to Z/S$.

\medskip\noindent
Suppose given objects $s_1 : X \to X_1 $ and $s_2 : X \to X_2$ in
$X/S$ and a morphism $a : X_1 \to X_2$ in $X/S$. Since $S$ is saturated,
we see that $a \in S$, see
Categories, Lemma \textnormal{[original label: \nolinkurl{categories-lemma-what-gets-inverted}]}.
By the argument of the previous paragraph we can complete
$s_1 : X \to X_1$ to an object
$(s_1, s'_1, s''_1) : (X, Y, Z, f, g, h) \to (X_1, Y_1, Z_1, f_1, g_1, h_1)$
in $\mathcal{I}$. Then we can repeat and find
$(a, b, c) : (X_1, Y_1, Z_1, f_1, g_1, h_1) \to (X_2, Y_2, Z_2, f_2, g_2, h_2)$
with $a, b, c \in S$ completing the given $a : X_1 \to X_2$.
But then $(a, b, c)$ is a morphism in $\mathcal{I}$.
In this way we conclude that the functor $\mathcal{I} \to X/S$ is
also surjective on arrows. Using rotation as above,
this implies the same thing is true for the functors
$\mathcal{I} \to Y/S$ and $\mathcal{I} \to Z/S$.

\medskip\noindent
The category $\mathcal{I}$ is nonempty as the identity provides an object.
This proves the condition (1) of the definition of a filtered category, see
Categories, Definition \textnormal{[original label: \nolinkurl{categories-definition-directed}]}.

\medskip\noindent
We check condition (2) of
Categories, Definition \textnormal{[original label: \nolinkurl{categories-definition-directed}]}
for the category $\mathcal{I}$. Suppose given objects
$(s_1, s'_1, s''_1) : (X, Y, Z, f, g, h) \to (X_1, Y_1, Z_1, f_1, g_1, h_1)$
and
$(s_2, s'_2, s''_2) : (X, Y, Z, f, g, h) \to (X_2, Y_2, Z_2, f_2, g_2, h_2)$
in $\mathcal{I}$. We want to find an object of $\mathcal{I}$
which is the target of an arrow from both
$(X_1, Y_1, Z_1, f_1, g_1, h_1)$ and $(X_2, Y_2, Z_2, f_2, g_2, h_2)$.
By Categories, Remark
\textnormal{[original label: \nolinkurl{categories-remark-left-localization-morphisms-colimit}]}
the categories $X/S$, $Y/S$, $Z/S$ are filtered.
Thus we can find $X \to X_3$ in $X/S$ and morphisms
$s : X_2 \to X_3$ and $a : X_1 \to X_3$. By the above we can find a morphism
$(s, s', s'') : (X_2, Y_2, Z_2, f_2, g_2, h_2) \to
(X_3, Y_3, Z_3, f_3, g_3, h_3)$ with $s', s'' \in S$.
After replacing $(X_2, Y_2, Z_2)$ by $(X_3, Y_3, Z_3)$ we may
assume that there exists a morphism $a : X_1 \to X_2$ in $X/S$.
Repeating the argument for $Y$ and $Z$ (by rotating as above)
we may assume there is a morphism
$a : X_1 \to X_2$ in $X/S$,
$b : Y_1 \to Y_2$ in $Y/S$, and
$c : Z_1 \to Z_2$ in $Z/S$.
However, these morphisms do not necessarily give rise to a morphism of
distinguished triangles. On the other hand, the necessary diagrams
do commute in $S^{-1}\mathcal{D}$. Hence we see (for example) that
there exists a morphism $s'_2 : Y_2 \to Y_3$ in $S$ such that
$s'_2 \circ f_2 \circ a = s'_2 \circ b \circ f_1$. Another replacement
of $(X_2, Y_2, Z_2)$ as above then gets us to the situation where
$f_2 \circ a = b \circ f_1$. Rotating and applying the same argument
two more times we see that we may assume $(a, b, c)$ is a morphism
of triangles. This proves condition (2).

\medskip\noindent
Next we check condition (3) of
Categories, Definition \textnormal{[original label: \nolinkurl{categories-definition-directed}]}.
Suppose $(s_1, s_1', s_1'') : (X, Y, Z) \to (X_1, Y_1, Z_1)$ and
$(s_2, s_2', s_2'') : (X, Y, Z) \to (X_2, Y_2, Z_2)$
are objects of $\mathcal{I}$, and suppose $(a, b, c), (a', b', c')$
are two morphisms between them. Since $a \circ s_1 = a' \circ s_1$
there exists a morphism $s_3 : X_2 \to X_3$ such that
$s_3 \circ a = s_3 \circ a'$. Using the surjectivity statement
we can complete this to a morphism of triangles
$(s_3, s_3', s_3'') : (X_2, Y_2, Z_2) \to (X_3, Y_3, Z_3)$
with $s_3, s_3', s_3'' \in S$. Thus
$(s_3 \circ s_2, s_3' \circ s_2', s_3'' \circ s_2'') :
(X, Y, Z) \to (X_3, Y_3, Z_3)$ is also an object of $\mathcal{I}$
and after composing the maps $(a, b, c), (a', b', c')$ with
$(s_3, s_3', s_3'')$ we obtain $a = a'$. By rotating we may do the
same to get $b = b'$ and $c = c'$.

\medskip\noindent
Finally, we check that $\mathcal{I} \to X/S$ is cofinal, see
Categories, Definition \textnormal{[original label: \nolinkurl{categories-definition-cofinal}]}.
The first condition is true as the functor is surjective.
Suppose that we have an object $s : X \to X'$ in $X/S$ and
two objects
$(s_1, s'_1, s''_1) : (X, Y, Z, f, g, h) \to (X_1, Y_1, Z_1, f_1, g_1, h_1)$
and
$(s_2, s'_2, s''_2) : (X, Y, Z, f, g, h) \to (X_2, Y_2, Z_2, f_2, g_2, h_2)$
in $\mathcal{I}$ as well as morphisms $t_1 : X' \to X_1$ and
$t_2 : X' \to X_2$ in $X/S$. By property (2) of $\mathcal{I}$
proved above we can find morphisms
$(s_3, s'_3, s''_3) : (X_1, Y_1, Z_1, f_1, g_1, h_1) \to
(X_3, Y_3, Z_3, f_3, g_3, h_3)$
and
$(s_4, s'_4, s''_4) : (X_2, Y_2, Z_2, f_2, g_2, h_2) \to
(X_3, Y_3, Z_3, f_3, g_3, h_3)$ in $\mathcal{I}$.
We would be done if the compositions
$X' \to X_1 \to X_3$ and $X' \to X_2 \to X_3$ were equal
(see displayed equation in 
Categories, Definition \textnormal{[original label: \nolinkurl{categories-definition-cofinal}]}).
If not, then, because $X/S$ is filtered, we can choose
a morphism $X_3 \to X_4$ in $X/S$ such that the compositions
$X' \to X_1 \to X_3 \to X_4$ and $X' \to X_2 \to X_3 \to X_4$ are equal.
Then we finally complete $X_3 \to X_4$ to a morphism
$(X_3, Y_3, Z_3) \to (X_4, Y_4, Z_4)$ in $\mathcal{I}$
and compose with that morphism to see that the result is true.
\end{proof}










\bibliographystyle{plain}
\bibliography{my}
\end{document}
